\documentclass[10pt,a4paper]{article}

\usepackage[margin=1.05in]{geometry}
\usepackage{setspace}
\usepackage{amsmath,amssymb,amsthm}
\usepackage{graphicx}
\usepackage{array}
\usepackage{caption}
\usepackage{enumitem}
\usepackage{titlesec}
\usepackage{microtype}
\usepackage{hyperref}
\usepackage{fancyhdr}

\hypersetup{colorlinks=true,linkcolor=black,urlcolor=black,citecolor=black}
\captionsetup{font=small,labelfont=bf,skip=4pt,justification=raggedright}

\titleformat{\section}{\normalsize\bfseries}{}{0em}{}
\titleformat{\subsection}{\small\bfseries}{}{0em}{}
\titleformat{\subsubsection}{\small\itshape}{}{0em}{}
\titlespacing*{\section}{0pt}{1.35em}{0.5em}
\titlespacing*{\subsection}{0pt}{0.95em}{0.35em}

\theoremstyle{plain}
\newtheorem{proposition}{Proposition}
\newtheorem{lemma}{Lemma}
\newtheorem{theorem}{Theorem}
\newtheorem{corollary}{Corollary}
\theoremstyle{definition}
\newtheorem{definition}{Definition}
\theoremstyle{remark}
\newtheorem{remark}{Remark}

\setstretch{1.14}
\setlist[itemize]{leftmargin=1.15em,itemsep=0.1em,topsep=0.2em}

\graphicspath{{drawings/}}

\pagestyle{fancy}
\fancyhf{}
\fancyhead[L]{\small\itshape Elements of Position}
\fancyhead[R]{\small\itshape IV.\ The Cut}
\fancyfoot[C]{\small\thepage}
\renewcommand{\headrulewidth}{0pt}

\begin{document}

\begin{center}
\vspace*{2.2em}
{\large Elements of Position}\\[1.6em]
{\LARGE\bfseries IV.\ The Cut}\\[2.2em]
{\normalsize Justin Erholtz}\\[0.35em]
{\small 12 September 2026}
\end{center}

\vspace{2.4em}

\begin{center}
{\large Every meeting sits on the line the unit already had.}
\end{center}

\vspace{1.8em}

\noindent
Parts I--III built a unit, a second walk, and the meetings of those walks.
This part asks what line those meetings are allowed to occupy.

The answer is not a new cut --- it is the only primary self-ratio the construction ever had: $1/2$.

Rate $0$ and rate not $0$ are two poles of one relation, not two worlds; scale is what you write in the title block.
Infinity is not a tick of the $n$-count, and a name is not a seat.

When the three drawings of one object --- plan, section, elevation --- are read together at finished stations, a later classical name may be used for that reading.
It does not get the first word.
The cut does.

\vspace{1.4em}

\noindent
\textit{A note.}
This part cites I, II, and III and stays inside the same questions.
Claims that can be proved here are proved by school geometry, or by a named identification after the drawings exist together.
The drawings are layers on the same architectural set, built by \texttt{Elements of Position - Drawings.py}.

\vspace{0.9em}
\noindent
\textit{A note on names used here.}
Names from I--III stand.
Cut --- the even split $1/2$, the only primary self-ratio of a finished segment.
Pause --- rate $0$ at one HERE.
Seat --- a finished comparison the construction can occupy from one own-origin.
Derived unit --- the reading of plan, section, and elevation together in the limit of finished stations.
Meeting of the derived unit --- a vanishing of a formula the stock can actually set to zero, from one HERE.

\newpage
\tableofcontents
\newpage

% =====================================================================
\section{What is already on the page}
% =====================================================================

I closed the unit, inverted the pair, named the diameter-$1$ close, and ran the $n$-count with two measures of one segment: axis length $n$, half-walk $n\pi/2$.
They lock at $n=2$, and the even cut is $1/2$ at every $n$.

II named $2$ from $1$, kept the rate at $1$, read irreducible interiors after they arrive, and laid $\Phi=1/\varphi$ on the same midline as a second process, not a second cut.

III watched the two processes as a system and solved the five-disk.
The first exact marriage sits on that same line, approximate dockings approach it, and the line does not move.

The missing sentence was not height.
It was what a vanishing is allowed to be on a trace that has one midpoint.

% =====================================================================
\section{Vacuum, trace, and rate as one scale}
% =====================================================================

Vacuum is the respect in which a relation may be named without being walked.
The inversion $x\mapsto 1/x$ can be written in one stroke, and infinity may be named.

Trace is the same relation inscribed as positions; with a finite speed of comparison there is a path-cost.
The cost vanishes where a pair coincides and grows where vacuum scales freely.

Naming is not arrival.
Collapse the two and a name is treated as a seat: privilege returns by the back door.

Rate is the other pole of that pair, not a second clock and not a second geometry.
Rate $0$ is pause: process off, no next tick, no new arrival.
What remains live is the unfinished expansion of the local circle already owned at that station.
Rate other than $0$ is trace: the $n$-count is moving.

The whole construction is scale-invariant, including rate: no unit is attached to the tempo.
One day between integer counts, one year, one millisecond --- the stick changes, the drawing does not.
The count is still $1,2,3,\ldots$, infinity is still not a tick, and no tempo reaches $\infty$ as a trace.

That sentence retires a hiding place: an extra zero cannot be waiting at completed infinity.
There is no such tick on this street.

Rate $0$ is not a second world in which a finished object lives with extra rooms; it is a pause of the same object, at one HERE, with one own-origin and one local circle.
Stopping does not mint a second centre.

\begin{figure}[ht]
\centering
\includegraphics[width=0.92\textwidth]{layer_d04_poles.png}
\caption{Two poles of one relation.
Rate $0$ is pause.
Any other rate is the $n$-count moving.}
\end{figure}

\begin{figure}[ht]
\centering
\includegraphics[width=0.92\textwidth]{layer_d02_here_there.png}
\caption{HERE and THERE are the same unit, other writing.}
\end{figure}

% =====================================================================
\section{Street, unit, THERE}
% =====================================================================

A street is an independent rigid process already owned by the closed unit: it carries its own ratios and its own infinity.
A unit is a finished relative comparison among streets: trace precedes measurement, measurement precedes unit.
The construction appoints neither a unit nor a vanishing.

THERE is not a mark waiting on the ruler --- it is elected when a comparison arrives.
A unit the construction can seat sits on the $1/2$ street; change the comparison and the unit changes, but the street does not.

A live half-walk $n\pi/2$ is the half-turn street measuring itself, not a derived comparison and not a seated meeting of a later unit.
III already refused that identification.
It is refused again here.

The closed unit already owns three drawings of one object.
Plan: the axis, the $n$-count, the pair $\{n,1/n\}$, the nested packages.
Section: the local circle of radius $n/2$, half-walk $n\pi/2$, first lock $\pi$ at $n=2$.
Elevation: inversion through $1$, look-back against what has already been built.

The derived unit of this part is the reading of those three drawings together in the limit of finished stations.
A meeting of that unit is a vanishing from one HERE.
Different later formulae are different routes across the same drawings; they elect different meetings, but they do not compete for the street.

% =====================================================================
\section{Complementary interiors}
% =====================================================================

Look-back from a pause is logarithm; live ticks are additive.
The invariant of inversion is $dx/x$; its integral along a finished arrival $n>1$ is $\log n$.

Let $\alpha$ be signed distance from $1/2$ on the unit, and set $\sigma=\tfrac12+\alpha$.
The inversion pair with $x=n^{\alpha}$ is $\{n^{\alpha},n^{-\alpha}\}$.
Multiply through by $n^{-1/2}$:
\[
\{n^{-\sigma},\,n^{-(1-\sigma)}\}.
\]

\begin{lemma}[Weights]
At a finished arrival $n>1$, the pair $\{n^{-\sigma},\,n^{-(1-\sigma)}\}$ has geometric mean $n^{-1/2}$ for every real $\sigma$.
The two ends are equal if and only if $\sigma=1/2$.
\end{lemma}

\begin{proof}
The product of the ends is $n^{-1}$, so the geometric mean is $n^{-1/2}$.
Equality of ends is $n^{-\sigma}=n^{-(1-\sigma)}$, hence $\sigma=1/2$.
Equivalently, $\{n^{\alpha},n^{-\alpha}\}$ meets if and only if $\alpha=0$.
\end{proof}

This is the first page of I, applied to complementary interiors: AM--GM on those two positive reals is minimized only at the even cut, and two leftovers are not two locks.
Axial leftover $\alpha=\sigma-\tfrac12$.
Circular leftover of the turn.
The midpoint lemma forbids a second primary ratio, but it does not forbid leftover of the circle.

\begin{figure}[ht]
\centering
\includegraphics[width=0.92\textwidth]{layer_d06_inversion.png}
\caption{Inversion through the unit.
The pair meets only at the lock.}
\end{figure}

% =====================================================================
\section{The only primary ratio, pulled back}
% =====================================================================

\begin{lemma}[Only primary self-ratio]
On any real segment there is a unique primary, scale-invariant, $\pm$-equal self-description of location: the midpoint ratio $1/2$.
\end{lemma}

It is primary: it does not come after another cut.
It is scale-invariant: it is a ratio of parts, not a length.
It is $\pm$-equal: the same amount from both ends.
Nested midpoints are that ratio asked again, not new cuts.
The similar cuts $\Phi$ and $\Phi^{2}$ are maximal similar copies, not a second primary self-description of location.
II already sat them on the same midline.

\begin{definition}[Permanent analytic cut]
A permanent analytic cut of the derived unit is a vertical line $\operatorname{Re}s=\sigma_0$ that functions as a primary, scale-invariant, $\pm$-equal self-description of location for that unit under the rate-$1$ identification of live scale and vacuum respect.
\end{definition}

\begin{lemma}[Only permanent analytic cut]
Under rate $1$, the only permanent analytic cut of the derived unit is $\operatorname{Re}s=1/2$.
\end{lemma}

\begin{proof}
Rate $1$ identifies the live geometry and the vacuum respect on one scale.
By the midpoint lemma the unique primary self-ratio on a segment is $1/2$.
The segment $[0,1]$ in the live geometry is the strip $0<\operatorname{Re}s<1$ in the vacuum writing of the same unit.
That pull-back is named here and not before.
The midpoint of $[0,1]$ is $1/2$.
Its image is the line $\operatorname{Re}s=1/2$.
Any other vertical line $\operatorname{Re}s=\sigma_0$ with $\sigma_0\neq 1/2$ would have to function as a primary self-description of location for that strip.
There is no such description other than the midpoint ratio.
\end{proof}

Uniqueness of the cut is not yet uniqueness of every vanishing.
A pair of names $\{\rho,1-\rho\}$ still has midpoint $1/2$, but the pairing does not collapse the pair.
A meeting is extra: the two ends meet.

% =====================================================================
\section{What the stock can set to zero}
% =====================================================================

A finished irreducible arrival $p$ is a station.
The data permitted there are the arrivals $q\le p$, the local circle of radius $p/2$, and look-back against what has already arrived.
A station is not a meeting, but a place at which one may pause.

On that station the dated product
\[
\zeta_p(s)=\prod_{q\le p}(1-q^{-s})^{-1}
\]
is a finite object.
No arrival beyond $p$ is present.
Each factor $(1-q^{-s})^{-1}$ is never zero wherever it is defined, so the dated product never vanishes.

The local-circle sample has no zeros in $0<\operatorname{Re}s<1$.
A midline polynomial $\tfrac12 s(s-1)$ accounts for the endpoints $0$ and $1$, and stops there.

Therefore a non-trivial vanishing is not a finite-depth fact.
It appears only when stations deepen, and only in a formula that is allowed to be zero.
In this stock there is one such formula: the dual sum of the complementary pair from the weights lemma.
The product cannot do it.
The circle cannot do it.
The sum of the two ends is what remains.

\begin{theorem}[Only vanishing mechanism]
In the strip $0<\operatorname{Re}s<1$, the only formula in the stock that can vanish non-trivially is the dual sum of the pair.
\end{theorem}

% =====================================================================
\section{A pause is one centre}
% =====================================================================

Write $\chi(s)$ for the ratio of local-circle samples at $1-s$ and $s$.
For $t\in\mathbb{R}$,
\[
\lvert\chi(\sigma+it)\rvert=1\quad\text{if and only if}\quad\sigma=\tfrac12.
\]
Off the lock, $\lvert\chi(\sigma+it)\rvert$ is the ratio of two radii, two own-origins.

A pause occupies one own-origin.
The $n$-count does not seat two centres at once.
Therefore a pause-reading has one circle, hence $\lvert\chi\rvert=1$, hence $\sigma=1/2$.

The off-lock dual sum, with $\chi$ compensating unequal ends so they can still add to zero, is not a pause of this construction --- it is two centres named at once.
Vacuum may name it, but trace cannot walk it, and a vanishing is an arrival: privilege does not return by the back door.

\begin{theorem}[One centre]
A vanishing written as a pause-reading has $\lvert\chi\rvert=1$, hence $\sigma=1/2$.
Combined with the weights lemma, the two ends meet only at half.
\end{theorem}

% =====================================================================
\section{Meetings sit on the cut}
% =====================================================================

\begin{lemma}[Meetings sit on the cut]
Every meeting of the derived unit that the construction can seat has real part $1/2$.
\end{lemma}

\begin{proof}
A meeting of the derived unit is a vanishing of a formula in the stock, from one HERE.
The only such formula in the strip is the dual sum of the pair.
A vanishing written as a pause-reading has one own-origin, hence $\lvert\chi\rvert=1$, hence $\sigma=1/2$.

Suppose an off-lock pair $\{\rho,1-\rho\}$ with $\operatorname{Re}\rho\neq 1/2$.
The pairing $s\leftrightarrow 1-s$ names both ends, but naming is vacuum, and a name is not a unit.
To seat the pair as a meeting the construction would have to occupy two own-origins at once, and the $n$-count does not.
Therefore the pair is not a seated meeting of this unit.

Therefore every meeting the derived unit can seat has real part $1/2$.
\end{proof}

A naming choice belongs here, ahead of the theorem rather than trailing after it.
The theorem below is a statement about the derived unit exactly as built in this stock: it proves nothing about any object outside it.
If that derived unit is identified with the classical limit obtained by reading the three drawings together through finished stations, the same theorem also reads as the classical statement that every non-trivial zero of that limit has real part $1/2$.
That identification is a name, not a derivation.
It is granted only once the drawings exist together; it is not an ingredient of I, II, or III, and no step of the proof below depends on it.

\begin{theorem}[Midline]
Every meeting of the derived unit has real part $1/2$.
\end{theorem}

\begin{proof}
A meeting of the derived unit is a seated vanishing from one HERE.
Every meeting the construction can seat has real part $1/2$.
An unseated name is not a meeting of the unit.
\end{proof}

The corollary is a statement about $1/2$: a meeting off the cut would be two centres named as one arrival, or a second primary self-description of location, and the geometry admits neither.

% =====================================================================
\section{What stands}
% =====================================================================

Forced:
the only primary self-ratio on a segment is $1/2$;
rate $0$ and rate not $0$ as poles of one relation;
no tempo reaches infinity as a trace;
a pause is one own-origin;
complementary interiors meet if and only if $\sigma=1/2$;
the dated product at a station never vanishes;
$\lvert\chi\rvert=1$ if and only if $\sigma=1/2$;
every seated meeting of the derived unit has real part $1/2$.

Named:
cut, pause, seat, derived unit, meeting of the derived unit, midline.

Not granted:
a second primary ratio;
two centres at one pause;
a name treated as a seat;
a live half-walk treated as a meeting of the derived unit;
a completed $n$-count as a first ingredient;
infinity as a last index.

The first sentence is still the sentence.
Where is $1$?
On the cut.
Walk there and you are HERE.
The meetings that can be seated sit on that same line, because there is no other primary ratio for them to occupy.

\end{document}
